\documentclass[11pt,a4paper]{article}

%================================================
% ENCODAGE ET LANGUE
%================================================
\usepackage[utf8]{inputenc}
\usepackage[T1]{fontenc}
\usepackage[french]{babel}

%================================================
% MISE EN PAGE
%================================================
\usepackage{geometry}
\geometry{margin=1.55cm}

%================================================
% PACKAGE GTOLI
%================================================
\usepackage{../styles/gtoli}

%================================================
% INFORMATIONS DU DOCUMENT
%================================================
\title{\textbf{Chimie}\\[0.2cm]
\large Titre du chapitre}

\author{GToli}
\date{\today}

%================================================
% DOCUMENT
%================================================
\begin{document}

\maketitle

%================================================
\section{Notion chimique concernée}
%================================================

\begin{cadreobjectif}
Dans ce chapitre, l’élève doit être capable de :
\begin{itemize}
    \item identifier les espèces chimiques en jeu ;
    \item utiliser correctement le vocabulaire chimique ;
    \item écrire ou interpréter une équation chimique ;
    \item appliquer une loi ou une relation chimique ;
    \item effectuer les calculs avec les unités adéquates ;
    \item interpréter un résultat dans le contexte chimique.
\end{itemize}
\end{cadreobjectif}

%================================================
\section{Situation de départ}
%================================================

\begin{cadrebleu}{Problématique chimique}
Cette section permet d’introduire la situation chimique étudiée.

\medskip

On précisera :
\begin{itemize}
    \item les substances ou espèces chimiques concernées ;
    \item le phénomène observé ;
    \item la question scientifique posée ;
    \item les hypothèses initiales éventuelles.
\end{itemize}
\end{cadrebleu}

%================================================
\section{Espèces chimiques et vocabulaire}
%================================================

\begin{cadretheorie}
\begin{itemize}
    \item \textbf{Espèce chimique :} entité chimique à préciser.
    \item \textbf{Réactif :} espèce consommée au cours d’une réaction.
    \item \textbf{Produit :} espèce formée au cours d’une réaction.
    \item \textbf{Coefficient stœchiométrique :} nombre placé devant une formule chimique dans une équation équilibrée.
    \item \textbf{Quantité de matière :} grandeur exprimée en mole.
\end{itemize}
\end{cadretheorie}

%================================================
\section{Tableau des grandeurs chimiques}
%================================================

\begin{cadregris}{Grandeurs, symboles et unités}
\renewcommand{\arraystretch}{1.35}
\begin{tabularx}{\textwidth}{>{\bfseries}p{3.2cm}p{2.2cm}p{2.5cm}X}
\toprule
Grandeur & Symbole & Unité & Signification \\
\midrule
Masse & $m$ & g & Masse d’un échantillon \\
Masse molaire & $M$ & g/mol & Masse d’une mole d’espèce chimique \\
Quantité de matière & $n$ & mol & Nombre de moles \\
Concentration molaire & $C$ & mol/L & Quantité de matière par litre de solution \\
Volume & $V$ & L & Volume de solution ou de gaz \\
\bottomrule
\end{tabularx}
\end{cadregris}

%================================================
\section{Rappel théorique}
%================================================

\begin{cadretheorie}
Cette section contient :
\begin{itemize}
    \item les définitions importantes ;
    \item les relations chimiques ;
    \item les formules ;
    \item les conventions de notation ;
    \item les conditions d’application ;
    \item les unités à respecter.
\end{itemize}

\medskip

\textbf{Point essentiel :} en chimie, il faut relier les espèces chimiques, les quantités de matière, les unités et le sens de la transformation.
\end{cadretheorie}

%================================================
\section{Relations fondamentales}
%================================================

\begin{cadrevert}{Formules utiles}
Relations fréquentes :

\[
n=\frac{m}{M}
\]

\[
C=\frac{n}{V}
\]

\[
\rho=\frac{m}{V}
\]

\[
pH=-\log[H_3O^+]
\]

\medskip

\textbf{Attention :} chaque formule doit être utilisée avec des unités cohérentes.
\end{cadrevert}

%================================================
\section{Équation chimique}
%================================================

\begin{cadremethode}
Pour exploiter une équation chimique :

\begin{enumerate}
    \item identifier les réactifs ;
    \item identifier les produits ;
    \item écrire les formules chimiques ;
    \item équilibrer l’équation ;
    \item lire les coefficients stœchiométriques ;
    \item établir les proportions molaires ;
    \item interpréter le sens de la réaction.
\end{enumerate}
\end{cadremethode}

\begin{cadregris}{Exemple d’équation}
\[
aA + bB \longrightarrow cC + dD
\]

où $a$, $b$, $c$ et $d$ sont les coefficients stœchiométriques.
\end{cadregris}

%================================================
\section{Tableau d’avancement}
%================================================

\begin{cadregris}{Modèle de tableau}
\renewcommand{\arraystretch}{1.35}
\begin{tabularx}{\textwidth}{>{\bfseries}p{3cm}X X X}
\toprule
État & Réactif 1 & Réactif 2 & Produit \\
\midrule
Initial & $n_{1,i}$ & $n_{2,i}$ & $0$ \\
Variation & $-a\xi$ & $-b\xi$ & $+c\xi$ \\
Final & $n_{1,i}-a\xi$ & $n_{2,i}-b\xi$ & $c\xi$ \\
\bottomrule
\end{tabularx}
\end{cadregris}

%================================================
\section{Protocole de résolution}
%================================================

\begin{cadremethode}
Pour résoudre un exercice de chimie :

\begin{enumerate}
    \item lire attentivement l’énoncé ;
    \item identifier les espèces chimiques ;
    \item relever les données numériques ;
    \item convertir les unités si nécessaire ;
    \item choisir la relation chimique adaptée ;
    \item écrire ou équilibrer l’équation si nécessaire ;
    \item travailler en quantité de matière lorsque c’est pertinent ;
    \item effectuer les calculs ;
    \item vérifier les unités ;
    \item interpréter chimiquement le résultat.
\end{enumerate}
\end{cadremethode}

%================================================
\section{Schéma ou représentation}
%================================================

\begin{cadregris}{Représentation simplifiée}
Insérer ici un schéma de montage expérimental, une représentation microscopique, un diagramme énergétique, une courbe cinétique ou un schéma de transformation.

\medskip

\begin{center}
\begin{tikzpicture}
\draw[thick] (0,0) rectangle (2,3);
\draw[fill=blue!10] (0.15,0.15) rectangle (1.85,1.7);
\node at (1,2.4) {solution};
\node at (1,0.9) {$A_{(aq)}$};
\draw[->,thick] (2.5,1.5) -- (4.2,1.5);
\node at (3.35,1.85) {réaction};
\draw[thick] (4.7,0) rectangle (6.7,3);
\draw[fill=green!10] (4.85,0.15) rectangle (6.55,1.7);
\node at (5.7,0.9) {$B_{(aq)}$};
\end{tikzpicture}
\end{center}
\end{cadregris}

%================================================
\section{Exemple modèle}
%================================================

\begin{cadreexemple}
\textbf{Énoncé :}

Insérer ici un exercice représentatif du chapitre.

\medskip

\textbf{Analyse chimique :}
\begin{itemize}
    \item Quelles espèces chimiques interviennent ?
    \item S’agit-il d’une transformation chimique, d’une dissolution, d’un dosage ou d’un équilibre ?
    \item Quelles grandeurs sont connues ?
    \item Quelle relation faut-il utiliser ?
\end{itemize}

\medskip

\textbf{Résolution :}

Présenter les étapes avec unités, conversions et interprétation.
\end{cadreexemple}

%================================================
\section{Expérience ou manipulation}
%================================================

\begin{cadreactivite}
Cette section peut contenir :
\begin{itemize}
    \item le matériel ;
    \item les consignes de sécurité ;
    \item le protocole expérimental ;
    \item le tableau de résultats ;
    \item l’exploitation des mesures ;
    \item la conclusion expérimentale.
\end{itemize}
\end{cadreactivite}

%================================================
\section{Sécurité}
%================================================

\begin{cadreattention}
En chimie, il faut toujours :
\begin{itemize}
    \item lire les pictogrammes de danger ;
    \item porter les protections nécessaires ;
    \item ne jamais goûter une substance ;
    \item éviter les mélanges non demandés ;
    \item respecter les consignes de manipulation ;
    \item éliminer les déchets selon les indications.
\end{itemize}
\end{cadreattention}

%================================================
\section{Exercice d’application}
%================================================

\begin{cadreenonce}
\textbf{Exercice 1 :}

Insérer ici un exercice d’application directe.
\end{cadreenonce}

\begin{cadreaide}
Pour commencer :
\begin{itemize}
    \item identifier les espèces chimiques ;
    \item noter les données ;
    \item convertir les unités ;
    \item choisir la formule ;
    \item vérifier le sens chimique du résultat.
\end{itemize}
\end{cadreaide}

%================================================
\section{Solution détaillée}
%================================================

\begin{cadresolution}
La solution doit contenir :
\begin{enumerate}
    \item les données ;
    \item les conversions ;
    \item la relation utilisée ;
    \item les calculs ;
    \item les unités ;
    \item l’interprétation chimique.
\end{enumerate}
\end{cadresolution}

%================================================
\section{Erreurs fréquentes}
%================================================

\begin{cadreattention}
Les erreurs fréquentes en chimie sont :
\begin{itemize}
    \item confondre masse et quantité de matière ;
    \item oublier de convertir les unités ;
    \item mal équilibrer une équation ;
    \item utiliser les masses au lieu des moles dans les proportions ;
    \item oublier les coefficients stœchiométriques ;
    \item interpréter un résultat sans vérifier son sens chimique.
\end{itemize}
\end{cadreattention}

%================================================
\section{Exercices progressifs}
%================================================

\begin{cadreactivite}
\begin{enumerate}
    \item identifier les espèces chimiques ;
    \item équilibrer une équation ;
    \item calculer une quantité de matière ;
    \item utiliser une concentration ;
    \item exploiter une équation chimique ;
    \item résoudre un exercice de synthèse.
\end{enumerate}
\end{cadreactivite}

%================================================
\section{Synthèse}
%================================================

\begin{cadresynthese}
À retenir :

\begin{itemize}
    \item espèces chimiques importantes : \dotfill
    \item relation principale : \dotfill
    \item unité essentielle : \dotfill
    \item équation chimique : \dotfill
    \item erreur à éviter : \dotfill
    \item interprétation chimique : \dotfill
\end{itemize}
\end{cadresynthese}

%================================================
\section{Auto-évaluation}
%================================================

\begin{cadregris}{Je vérifie ma maîtrise}
\begin{itemize}
    \item[$\square$] Je connais le vocabulaire chimique.
    \item[$\square$] Je sais identifier réactifs et produits.
    \item[$\square$] Je sais équilibrer une équation.
    \item[$\square$] Je sais utiliser les unités correctes.
    \item[$\square$] Je sais choisir la bonne formule.
    \item[$\square$] Je sais interpréter le résultat.
\end{itemize}
\end{cadregris}

\end{document}